\documentclass[UTF8,a4paper,11pt]{ctexart}

\usepackage{amsmath,amssymb,amsthm}
\usepackage{geometry}
\usepackage{hyperref}
\usepackage{enumitem}

\geometry{margin=1in}
\hypersetup{
  colorlinks=true,
  linkcolor=blue,
  citecolor=blue,
  urlcolor=blue
}

\newtheorem{theorem}{定理}
\newtheorem{assumption}{假设}
\newtheorem{remark}{说明}

\title{QR Box-Dual FedPDHG 的收敛阶数说明}
\author{FED-QR 项目}
\date{\today}

\begin{document}

\maketitle

\begin{abstract}
本文给出 QR box-dual FedPDHG 方法的一个可用于报告或论文草稿的中文收敛性说明。
核心结论是：在全客户端、全样本、确定性更新下，QR box-dual PDHG 可以得到
$\mathcal{O}(1/T)$ 的遍历 primal-dual gap 收敛阶数；在随机联邦设定下，即存在
客户端随机抽样、用户 mini-batch 抽样以及 cached direction 的 stale 误差时，
不用方差缩减也可以证明 $\mathcal{O}(1/\sqrt{T})$ 阶的随机收敛率。固定步长时，
方法通常收敛到一个由随机方差和 stale bias 决定的误差邻域。方差缩减不是收敛
证明的必要条件，它主要降低随机方差常数并改善稳定性。
\end{abstract}

\section{问题形式与对偶盒结构}

考虑带惩罚项的分位数回归问题
\begin{equation}
    \min_{\beta\in\mathbb{R}^p}
    F(\beta)
    :=
    \frac{1}{n}
    \sum_{i=1}^n
    \rho_\tau(y_i-x_i^\top\beta)
    +
    r(\beta),
    \label{eq:primal}
\end{equation}
其中
\[
    \rho_\tau(u)=u(\tau-\mathbf{1}\{u<0\})
\]
是分位数回归的 check loss，$\tau\in(0,1)$ 是分位数水平，$r(\beta)$ 是惩罚项。
在理论部分，我们先假设 $r$ 是 proper、closed、convex 的函数；如果使用
MCP/SCAD 等非凸惩罚，则严格全局收敛结论需要改成 stationary-point 类型结论。

check loss 的关键结构是 Fenchel 对偶盒表示：
\begin{equation}
    \rho_\tau(u)
    =
    \max_{v\in[\tau-1,\tau]} vu .
\end{equation}
因此，原问题 \eqref{eq:primal} 可以写成 convex-concave saddle point 问题：
\begin{equation}
    \min_{\beta\in\mathbb{R}^p}
    \max_{v\in\mathcal{V}}
    \mathcal{L}(\beta,v)
    :=
    r(\beta)
    +
    \frac{1}{n}
    v^\top (y-X\beta),
    \qquad
    \mathcal{V}:=[\tau-1,\tau]^n .
    \label{eq:saddle}
\end{equation}
QR box-dual 方法的核心正是利用这个盒约束对偶变量：
\[
    v_i\in[\tau-1,\tau].
\]
算法每轮对 $v_i$ 做盒投影更新，并对 $\beta$ 做 primal proximal update。

\section{确定性全量版本：$\mathcal{O}(1/T)$}

首先考虑最理想的确定性版本：

\begin{itemize}[leftmargin=2em]
    \item 所有客户端每轮都参与；
    \item 每个客户端每轮都刷新全部本地样本的 dual variables；
    \item 不存在客户端抽样随机性；
    \item 不存在用户 mini-batch 随机性；
    \item 不存在 stale cached direction。
\end{itemize}

\begin{assumption}[确定性 PDHG 条件]
\label{assump:det}
惩罚项 $r$ 是 proper、closed、convex 的函数。矩阵 $X$ 的谱范数有限。
步长 $\eta,\sigma>0$ 满足
\begin{equation}
    \eta\sigma \frac{\|X\|_2^2}{n^2}<1 .
    \label{eq:stepsize}
\end{equation}
\end{assumption}

\begin{theorem}[确定性 QR box-dual PDHG 的遍历收敛阶数]
\label{thm:det}
在假设 \ref{assump:det} 下，令 $(\beta^t,v^t)$ 为全量确定性 QR box-dual PDHG
生成的迭代序列。定义遍历平均
\[
    \bar{\beta}^T
    =
    \frac{1}{T}\sum_{t=1}^T\beta^t,
    \qquad
    \bar{v}^T
    =
    \frac{1}{T}\sum_{t=1}^Tv^t.
\]
若 $(\beta^\star,v^\star)$ 是 saddle point，则有
\begin{equation}
    \mathcal{G}(\bar{\beta}^T,\bar{v}^T)
    :=
    \mathcal{L}(\bar{\beta}^T,v^\star)
    -
    \mathcal{L}(\beta^\star,\bar{v}^T)
    \le
    \frac{1}{T}
    \left(
        \frac{\|\beta^\star-\beta^0\|^2}{2\eta}
        +
        \frac{\|v^\star-v^0\|^2}{2\sigma}
    \right).
    \label{eq:det-rate}
\end{equation}
因此，全量确定性 QR box-dual PDHG 的遍历 primal-dual gap 收敛阶数为
\[
    \mathcal{O}(1/T).
\]
\end{theorem}

\begin{proof}[证明思路]
将 \eqref{eq:saddle} 写成标准 PDHG 形式：
\[
    r(\beta)+\langle K\beta,v\rangle-h^\ast(v),
    \qquad
    K=-X/n,
\]
其中
\[
    h^\ast(v)
    =
    -v^\top y/n+\iota_{\mathcal{V}}(v),
\]
$\iota_{\mathcal{V}}$ 是集合 $\mathcal{V}$ 的 indicator function。
dual update 正是 $h^\ast$ 的 proximal step，也就是带 affine shift 后投影到
$[\tau-1,\tau]^n$；primal update 是 $r$ 的 proximal step。

由标准 PDHG 一步下降不等式可得，对任意 $(\beta,v)$，
\begin{align}
    \mathcal{L}(\beta^{t+1},v)
    -
    \mathcal{L}(\beta,v^{t+1})
    &\le
    \frac{1}{2\eta}
    \left(
        \|\beta-\beta^t\|^2
        -
        \|\beta-\beta^{t+1}\|^2
    \right)
    \nonumber\\
    &\quad+
    \frac{1}{2\sigma}
    \left(
        \|v-v^t\|^2
        -
        \|v-v^{t+1}\|^2
    \right).
    \label{eq:one-step}
\end{align}
对 $t=0,\ldots,T-1$ 求和后，右侧平方距离项 telescoping。再令
$(\beta,v)=(\beta^\star,v^\star)$，并利用 $\mathcal{L}$ 对 $\beta$ 的凸性和
对 $v$ 的凹性，即可得到 \eqref{eq:det-rate}。
\end{proof}

\section{随机联邦版本：不用方差缩减时的 $\mathcal{O}(1/\sqrt{T})$}

真实的联邦 QR box-dual 方法不是确定性全量更新，而是包含：

\begin{itemize}[leftmargin=2em]
    \item 每轮只抽取一部分客户端；
    \item 被抽中的客户端只用 mini-batch 用户更新 dual variables；
    \item 未被抽中的客户端使用 cached direction；
    \item cached direction 会产生 stale bias；
    \item stochastic direction 存在随机方差。
\end{itemize}

令 $\widehat{g}^t$ 表示服务器第 $t$ 轮实际使用的随机 cached primal direction，
$g^t$ 表示理想全量方向。

\begin{assumption}[随机联邦条件]
\label{assump:stoch}
存在有限常数 $R,V,\Sigma,D$，使得：
\begin{enumerate}[leftmargin=2em]
    \item 协变量有界：$\|x_i\|\le R$；
    \item dual variable 天然有界：
    \[
        v_i^t\in[\tau-1,\tau],
        \qquad
        |v_i^t|\le V:=\max\{\tau,1-\tau\};
    \]
    \item stochastic direction 方差有界：
    \begin{equation}
        \mathbb{E}
        \left[
        \left.
        \|\widehat{g}^t-\mathbb{E}(\widehat{g}^t\mid\mathcal{F}_t)\|^2
        \right|
        \mathcal{F}_t
        \right]
        \le
        \Sigma^2;
        \label{eq:variance}
    \end{equation}
    \item stale cache bias 满足
    \begin{equation}
        \left\|
        \mathbb{E}(\widehat{g}^t\mid\mathcal{F}_t)-g^t
        \right\|
        \le
        \Delta_t,
        \qquad
        \frac{1}{T}
        \sum_{t=0}^{T-1}
        \Delta_t
        \le
        \bar{\Delta}_T;
        \label{eq:bias}
    \end{equation}
    \item 迭代点停留在有界 level set 中，其直径不超过 $D$。
\end{enumerate}
\end{assumption}

\begin{theorem}[随机联邦 QR box-dual 的收敛阶数]
\label{thm:stoch}
在假设 \ref{assump:stoch} 下，令随机联邦 QR box-dual 使用
\[
    \eta_t=\sigma_t=\gamma/\sqrt{T}
\]
形式的步长，其中 $\gamma$ 足够小以满足 PDHG 稳定条件。则其遍历 primal-dual gap
满足
\begin{equation}
    \mathbb{E}
    \left[
    \mathcal{G}(\bar{\beta}^T,\bar{v}^T)
    \right]
    \le
    \frac{C_1}{\sqrt{T}}
    +
    C_2\bar{\Delta}_T,
    \label{eq:stoch-rate}
\end{equation}
其中 $C_1$ 与初始距离、$\gamma$、随机方差上界 $\Sigma^2$ 有关，$C_2$ 与有界
level set 的直径有关。若平均 stale bias 满足
\[
    \bar{\Delta}_T=\mathcal{O}(T^{-1/2}),
\]
则
\[
    \mathbb{E}
    \left[
    \mathcal{G}(\bar{\beta}^T,\bar{v}^T)
    \right]
    =
    \mathcal{O}(T^{-1/2}).
\]
\end{theorem}

\begin{proof}[证明思路]
随机情形的证明是在确定性 PDHG 一步不等式的基础上加入随机误差项。
因为实际使用的是 $\widehat{g}^t$ 而不是 $g^t$，所以多出
\[
    \left\langle
    \widehat{g}^t-g^t,\,
    \beta^{t+1}-\beta
    \right\rangle .
\]
将误差分解为两部分：
\[
    \widehat{g}^t-g^t
    =
    \underbrace{
    \widehat{g}^t
    -
    \mathbb{E}(\widehat{g}^t\mid\mathcal{F}_t)
    }_{\text{martingale noise}}
    +
    \underbrace{
    \mathbb{E}(\widehat{g}^t\mid\mathcal{F}_t)
    -
    g^t
    }_{\text{stale cache bias}}.
\]
第一项是零均值随机噪声，由方差上界 \eqref{eq:variance} 控制；第二项由
\eqref{eq:bias} 控制。

对一步不等式求和后，确定性部分仍然 telescoping；随机噪声部分通过 Young
不等式产生
\[
    \sum_{t=0}^{T-1}
    \eta_t\Sigma^2
\]
量级的项；stale bias 产生
\[
    D\sum_{t=0}^{T-1}\Delta_t
\]
量级的项。两边除以 $\sum_{t=0}^{T-1}\eta_t$ 后得到
\[
    \frac{C}{\sum_{t=0}^{T-1}\eta_t}
    +
    \frac{\sum_{t=0}^{T-1}\eta_t^2\Sigma^2}
    {\sum_{t=0}^{T-1}\eta_t}
    +
    C_2\bar{\Delta}_T.
\]
取 $\eta_t=\Theta(T^{-1/2})$，即可得到
\eqref{eq:stoch-rate}。
\end{proof}

\section{固定步长时的误差邻域}

实际实验中我们常使用固定步长。固定步长下，随机误差不会完全消失，因此更合理的
结论是收敛到一个误差邻域：
\begin{equation}
    \mathbb{E}
    \left[
    \mathcal{G}(\bar{\beta}^T,\bar{v}^T)
    \right]
    \le
    \frac{C_3}{\eta T}
    +
    C_4\eta\Sigma^2
    +
    C_5\bar{\Delta}_T .
    \label{eq:fixed}
\end{equation}
其中第一项是 $\mathcal{O}(1/T)$ 的 transient term，后两项分别来自随机方差和
stale cache bias。因此固定步长不是严格收敛到零 gap，而是收敛到
variance/staleness-dependent neighborhood。

\section{stale、robust、adaptive 对理论的影响}

staleness-aware 和 robust/adaptive 方法在服务器聚合时使用客户端权重 $w_j^t$。
如果这些权重满足
\[
    0<w_{\min}\le w_j^t\le w_{\max}<1,
    \qquad
    \sum_{j=1}^m w_j^t=1,
\]
则上述随机收敛阶数不发生本质变化，只是常数和优化目标发生变化。

\paragraph{sample weighting.}
如果
\[
    w_j = n_j/n,
\]
则优化的是原始 sample-average empirical risk。

\paragraph{uniform/custom weighting.}
如果使用 uniform 或 custom client weighting，则理论对应的目标变为 client-weighted
分位数回归：
\begin{equation}
    F_w(\beta)
    =
    \sum_{j=1}^m
    w_j
    \frac{1}{n_j}
    \sum_{i\in\mathcal{D}_j}
    \rho_\tau(y_i-x_i^\top\beta)
    +
    r(\beta).
\end{equation}
这解释了为什么 robust 方法可能牺牲 global sample-average objective，却改善
worst-client loss。

\paragraph{adaptive weighting.}
adaptive weighting 可以看成优化一个缓慢变化的 client-weighted saddle problem。
如果权重有界且总变化量有限，例如
\[
    \sum_{t=0}^{T-1}
    \|w^{t+1}-w^t\|_1
    <
    \infty,
\]
则额外项可以作为 lower-order variation error 加入收敛界中。

\section{方差缩减的作用}

方差缩减不是证明收敛阶数的必要条件。不使用方差缩减时，随机联邦版本已经可以在
上述假设下得到 $\mathcal{O}(1/\sqrt{T})$ 阶。

control-variate VR 的作用是降低随机方差常数。若校正后的方向
$\widehat{g}_{\mathrm{vr}}^t$ 满足
\[
    \mathbb{E}
    \left[
    \left.
    \|\widehat{g}_{\mathrm{vr}}^t
    -
    \mathbb{E}(\widehat{g}_{\mathrm{vr}}^t\mid\mathcal{F}_t)\|^2
    \right|
    \mathcal{F}_t
    \right]
    \le
    \Sigma_{\mathrm{vr}}^2
    \le
    \Sigma^2,
\]
则定理 \ref{thm:stoch} 中的 $\Sigma^2$ 可以替换为更小的
$\Sigma_{\mathrm{vr}}^2$。因此 VR 改善的是常数，而不是在一般 nonsmooth
随机联邦设定下直接把阶数从 $\mathcal{O}(1/\sqrt{T})$ 变成
$\mathcal{O}(1/T)$。

\section{总结}

综上，理论结论可以概括为：
\begin{itemize}[leftmargin=2em]
    \item 全量确定性 QR box-dual PDHG：遍历 primal-dual gap 为
    $\mathcal{O}(1/T)$；
    \item 随机联邦 QR box-dual，不用方差缩减：在 bounded variance 和 bounded
    stale bias 条件下为 $\mathcal{O}(1/\sqrt{T})$；
    \item 固定步长随机版本：收敛到由随机方差和 stale bias 决定的误差邻域；
    \item robust/adaptive weighting 不改变基本阶数，但改变目标和常数；
    \item 方差缩减降低 variance constant，主要改善稳定性和公平性，不是主收敛
    证明的必要条件。
\end{itemize}

\end{document}
